\chapter*{Theoretical Foundations of AR-Based Games}
\begin{sloppy}

\section*{Executive summary}
AR-based games are not simply mobile games with a camera overlay. Their theoretical novelty lies in the fact that cognition, attention, action, and reward are distributed across two coupled environments at once: the simulated rule system and the player's actual surroundings. In the classic AR definition, the system combines real and virtual content, operates interactively in real time, and registers content in 3D; in game terms, this means the physical world becomes part of both the representation layer and the rule layer. The literature therefore treats successful AR games as hybrid systems whose design must cohere across psychology, spatial computing, and game mechanics.

The most useful theoretical lenses cluster into three groups. First, embodied cognition and situated learning explain why AR can make tasks feel more meaningful and memorable: players perceive, move, and decide in context, often off-loading memory and navigation onto the environment itself. Second, flow and self-determination theory explain why some AR play loops feel compelling: challenge-skill balance, clear feedback, autonomy, competence, and relatedness all predict enjoyment and persistence. Third, persuasive technology and gamification theory explain why AR games can reshape mobility and routines, but also why poorly designed systems drift into coercion, distraction, or ethically dubious ``always-on'' prompting.

Technically, AR games depend on a stack that conventional mobile games do not need to solve nearly as deeply: sensor fusion, visual-inertial odometry, SLAM, anchors, scene understanding, occlusion, and persistence. These are not just engineering details. They determine what counts as a stable game object, whether multiple players can inhabit the ``same'' augmented world, whether virtual objects feel embedded rather than pasted on top, and whether a location-based rule remains legible across sessions. The design space of AR games is therefore constrained and enabled by tracking quality, latency, drift, environmental texture, lighting, network availability, and privacy architecture.

Compared with conventional mobile games, AR games add physicality, context-awareness, place-dependent progression, and public-space risk. They also add design obligations: safety messaging, accessibility beyond visual overlays, responsible location-data practices, resilience to drift and occlusion failure, and fairness across geographies with unequal densities of points of interest. Research on the geography and safety of large-scale location-based play shows that these issues are not peripheral; they are part of the core UX and ethics of AR game design.

\section*{Theoretical frameworks}
Embodied cognition is foundational because AR games place cognition inside active sensorimotor coupling with the environment. Wilson's classic review distinguishes several claims especially relevant to AR play: cognition is situated, time-pressured, off-loads work onto the environment, and is ultimately for action. In AR games, those claims become design obligations. Players navigate with bodily motion, orient around landmarks, use the world as memory support, and solve tasks through movement rather than abstract menu selection alone. AR thus turns ``context'' from a backdrop into a computational partner.

Situated learning provides the complementary social and pedagogical account. Lave and Wenger's notion of legitimate peripheral participation frames learning as participation in authentic activity and communities of practice, not the transfer of detached facts. AR learning games and place-based ARGs operationalize this by tying clues, roles, and consequences to real sites, making the environment itself part of the epistemic frame. Reviews of mobile AR learning repeatedly describe its value as embodied, active, contextualized, place-based, and authentic, while studies such as Environmental Detectives and Mad City Mystery show how role-play and physical place can support inquiry and argumentation rather than only content recall.

Flow theory explains why AR play is engaging when it works well and frustrating when it does not. Csikszentmihalyi's model centers on the balance between challenge and skill: if the task is too hard relative to skill, anxiety emerges; if too easy, boredom follows; sustained enjoyment appears when both challenge and skill are high and closely matched. In AR games, the challenge-skill relation includes not only combat or puzzle difficulty but also locomotion effort, device handling, spatial interpretation, and environmental noise. That is why the same mechanic can feel elegant in a clear, stable setting and exhausting in a cluttered or poorly tracked one.

Self-determination theory gives the most robust motivational vocabulary for AR game design. Ryan and Deci identify autonomy, competence, and relatedness as basic psychological needs. Their game-specific work shows that need satisfaction during play predicts enjoyment, immersion, preference for future play, and short-term well-being; the PENS instrument was created precisely to measure this in games. For AR, the implications are direct: autonomy is supported by meaningful route and task choice, competence by readable feedback and calibratable challenge, and relatedness by synchronous or asynchronous co-presence around shared places and objects.

Persuasive technology and gamification theory explain the behavioral power of AR systems beyond ``fun.'' Fogg's behavior model argues that behavior occurs when motivation, ability, and a prompt converge. Deterding and colleagues define gamification as the use of game design elements in non-game contexts, and Hamari, Koivisto, and Sarsa's review concludes that gamification often has positive effects but that results are strongly context-dependent. AR-based games blur the boundary between games and gamified behavior systems because they can turn walking, visiting, scanning, and social coordination into rewarded acts. This is exactly why they can promote exercise and exploration, and exactly why they can also produce compulsion, distraction, or exploitative urgency if prompts and rewards are not carefully bounded.

\section*{Technical interaction models}
At the tracking layer, modern AR games rely on sensor fusion between camera data and inertial measurements. Official documentation from Google states that ARCore combines feature tracking from the camera with IMU data to estimate device pose, while its fundamentals page explicitly describes this as SLAM. Long before current mobile SDKs, core AR tracking literature established the same logic: sensor fusion is needed because AR pose estimation demands both high accuracy and low latency; MonoSLAM, PTAM, and later ORB-SLAM moved the field from prepared, marker-based settings toward robust real-time mapping in unprepared scenes; and later benchmarking work showed how commercial visual-inertial systems such as ARCore and ARKit compare in authentic indoor and outdoor conditions.

Spatial mapping matters because game logic needs a model of ``where things can be.'' ARCore's Depth API generates depth maps that support object occlusion and more realistic interaction; ARKit's People Occlusion similarly allows content to pass behind and in front of people. These features are not cosmetic. They change whether an object reads as inhabiting the same space as the player, which in turn affects presence, affordance legibility, and trust in the simulation. Put analytically, occlusion and scene understanding produce world-consistency: the sense that virtual rules obey the same spatial constraints as physical ones. That is one reason immersion-related reviews repeatedly link realism, interactivity, and control to presence and engagement in AR games.

Anchoring and persistence turn a momentary overlay into a game world. ARCore's anchor guidance notes that objects should remain close to anchors to avoid instability, while Cloud Anchors are designed specifically to persist and share content across users and sessions. Apple and Niantic provide parallel capabilities at different scales: Apple through world tracking and scene reconstruction features, and Niantic through VPS, persistent anchors, and colocalization. Niantic's documentation is especially useful for games because it frames precise localization, shared anchors, and stable alignment as prerequisites for persistent and multi-user AR tied to mapped real-world sites.

\section*{Mechanics at the boundary of digital and physical}
The distinctive mechanics of AR games are those that recruit real-world properties into rules. Location-based mechanics are the clearest example. In Ingress, the official framing is ``discover, strategize, collaborate'' across portals distributed through the built environment; in Pokémon GO, PokéStops and Gyms supply items and encounters; and recent studies of location-based play identify 11 distinct dynamics through which point-of-interest placement, habitat, and surroundings incentivize travel to specific areas. In other words, the map is not just scenery; it is a live balance sheet of effort, opportunity, scarcity, and social encounter.

Object affordances in AR are more literal than in conventional mobile games because virtual objects inherit constraints from physical surfaces and sight lines. Minecraft Earth's launch materials made this explicit: a build can be placed on a picnic table or patch of grass, turning everyday surfaces into build sites and adventures into local events. Peridot pushes the same logic toward creature interaction by requiring the camera and framing itself as an AR-first experience in which the character reacts to the real world. These mechanics succeed when the system can reliably infer where objects ``belong'' and when the fiction respects what the space actually affords.

Environmental constraints are equally important. Weather, darkness, crowd density, weak texture, GPS error, internet instability, and unsafe movement paths all shape what counts as a fair or feasible challenge. Niantic's VPS documentation explicitly notes that moving too fast, darkness or fog, dirty lenses, and poor connectivity degrade localization quality. That means designers cannot treat the world as a neutral stage. They have to design for variable observability, partial failure, and public-space unpredictability.

\section*{Differences from conventional mobile games}
The main difference between AR-based games and ordinary mobile games is not graphical style but the locus of interaction. Conventional mobile games typically treat the screen as the entire primary play space. AR games distribute play across screen, body, place, and social surroundings. As a result, UX problems that are secondary in conventional games become first-order design variables in AR: orientation, environmental legibility, safe movement, weather, local network quality, landmark density, privacy disclosure, and re-entry after interruption.

\section*{Design patterns and case studies}
The field's evolution is easiest to understand as a movement from proof-of-concept spatial overlays to mature hybrid systems where physical space is a first-class rule substrate. Early research prototypes such as ARQuake and Can You See Me Now? established the bodily and locative premise. Educational systems such as Environmental Detectives and Mad City Mystery deepened the situated-learning argument. Mainstream mobile titles then industrialized the category with POI databases, public events, and live-ops economies, while newer ACM and mixed-reality work has shifted attention toward shared-world AR, design patterns, low-code toolchains, and cooperative prop-based play.

\end{sloppy}
